\textbf{Phần 1. Trắc nghiệm lựa chọn}
\setcounterpageref{ex}{0}

\begin{ex} %[2H5N25]
Người ta thường dùng các chữ cái in hoa nào để đặt tên cho điểm?
\choice
{\True $A, B, C$}
{$a, b, c$}
{$1, 2, 3$}
{$x, y, z$}
\loigiai{
Đặt tên cho điểm dùng chữ in hoa $A, B, C, \dots$\\
\textbf{Đáp án của câu hỏi là A}
}
\end{ex}

\begin{ex} %[2H5N25]
Đoạn thẳng $MN$ được giới hạn bởi mấy điểm?
\choice
{$1$ điểm}
{\True $2$ điểm ($M$ và $N$)}
{$3$ điểm}
{Không có điểm nào}
\loigiai{
Đoạn thẳng $MN$ giới hạn bởi hai đầu mút $M$ và $N$.\\
\textbf{Đáp án của câu hỏi là B}
}
\end{ex}

\begin{ex} %[2H5H25]
Ba điểm cùng nằm trên một đường thẳng được gọi là
\choice
{Ba điểm không thẳng hàng}
{Ba điểm cắt nhau}
{Ba điểm song song}
{\True Ba điểm thẳng hàng}
\loigiai{
Ba điểm cùng nằm trên 1 đường thẳng là ba điểm thẳng hàng.\\
\textbf{Đáp án của câu hỏi là D}
}
\end{ex}

\begin{ex} %[2H5H25]
Hình ảnh sợi dây điện chùng xuống cho ta hình ảnh của
\choice
{Đoạn thẳng}
{Đường thẳng}
{\True Đường cong}
{Tia số}
\loigiai{
Sợi dây chùng cho ta hình ảnh đường cong.\\
\textbf{Đáp án của câu hỏi là C}
}
\end{ex}

\begin{ex} %[2H5H25]
Đoạn thẳng $AB$ dài $8\text{ cm}$, đoạn thẳng $CD$ dài $12\text{ cm}$. Đoạn thẳng $CD$ dài hơn đoạn thẳng $AB$ là
\choice
{\True $4\text{ cm}$}
{$3\text{ cm}$}
{$5\text{ cm}$}
{$6\text{ cm}$}
\loigiai{
$12 - 8 = 4\text{ cm}$.\\
\textbf{Đáp án của câu hỏi là A}
}
\end{ex}

\begin{ex} %[2H5H25]
Cho điểm $B$ nằm giữa hai điểm $A$ và $C$ thẳng hàng. Biết $AB = 5\text{ cm}, BC = 4\text{ cm}$. Độ dài đoạn thẳng $AC$ là
\choice
{$8\text{ cm}$}
{\True $9\text{ cm}$}
{$10\text{ cm}$}
{$1\text{ cm}$}
\loigiai{
$AC = AB + BC = 5 + 4 = 9\text{ cm}$.\\
\textbf{Đáp án của câu hỏi là B}
}
\end{ex}

\begin{ex} %[2H5V25]
Qua hai điểm phân biệt $A$ và $B$, ta vẽ được bao nhiêu đường thẳng?
\choice
{$2$ đường thẳng}
{$3$ đường thẳng}
{\True Chỉ $1$ đường thẳng duy nhất}
{Vô số đường thẳng}
\loigiai{
Qua 2 điểm chỉ vẽ được duy nhất 1 đường thẳng.\\
\textbf{Đáp án của câu hỏi là C}
}
\end{ex}

\begin{ex} %[2H5V25]
Cho 3 điểm thẳng hàng $M, N, P$ theo thứ tự đó. Biết $MP = 15\text{ cm}, MN = 7\text{ cm}$. Độ dài đoạn thẳng $NP$ là
\choice
{\True $8\text{ cm}$}
{$7\text{ cm}$}
{$9\text{ cm}$}
{$22\text{ cm}$}
\loigiai{
$NP = 15 - 7 = 8\text{ cm}$.\\
\textbf{Đáp án của câu hỏi là A}
}
\end{ex}

\begin{ex} %[2H5V25]
Một đoạn thẳng $AB$ dài $14\text{ cm}$. Điểm $M$ là điểm chính giữa (ở giữa) của đoạn thẳng $AB$. Độ dài đoạn thẳng $AM$ là
\choice
{$6\text{ cm}$}
{$8\text{ cm}$}
{$9\text{ cm}$}
{\True $7\text{ cm}$}
\loigiai{
$14 : 2 = 7\text{ cm}$.\\
\textbf{Đáp án của câu hỏi là D}
}
\end{ex}

\begin{ex} %[2H5V25]
Cho 4 điểm $A, B, C, D$ cùng nằm trên một đường thẳng. Có tất cả bao nhiêu đoạn thẳng nối từ các điểm đó?
\choice
{$4$ đoạn thẳng}
{\True $6$ đoạn thẳng ($AB, AC, AD, BC, BD, CD$)}
{$5$ đoạn thẳng}
{$8$ đoạn thẳng}
\loigiai{
Các đoạn: $AB, AC, AD, BC, BD, CD$ (gồm 6 đoạn thẳng).\\
\textbf{Đáp án của câu hỏi là B}
}
\end{ex}

\begin{ex} %[2H5V25]
Đoạn thẳng $MN$ dài $20\text{ cm}$, đoạn thẳng $PQ$ ngắn hơn đoạn thẳng $MN$ là $6\text{ cm}$. Độ dài đoạn thẳng $PQ$ là
\choice
{$13\text{ cm}$}
{$15\text{ cm}$}
{\True $14\text{ cm}$}
{$26\text{ cm}$}
\loigiai{
$20 - 6 = 14\text{ cm}$.\\
\textbf{Đáp án của câu hỏi là C}
}
\end{ex}

\begin{ex} %[2H5C25]
Cho 3 điểm phân biệt không thẳng hàng $A, B, C$. Nối từng cặp điểm lại với nhau ta được hình gì?
\choice
{Đoạn thẳng dài}
{Hình tứ giác}
{Đường thẳng}
{\True Hình tam giác}
\loigiai{
Nối 3 điểm không thẳng hàng tạo thành hình tam giác.\\
\textbf{Đáp án của câu hỏi là D}
}
\end{ex}

\textbf{Phần 2. Trắc nghiệm đúng sai}
\setcounterpageref{ex}{0}

\begin{ex} %[2H5H25]
Xét tính đúng sai của các khẳng định hình học:
\choiceTF
{\True Đoạn thẳng có độ dài đo được cụ thể.}
{\True Đường thẳng kéo dài mãi về hai phía.}
{Ba điểm bất kì luôn luôn thẳng hàng.}
{\True Ba điểm chỉ thẳng hàng khi chúng cùng nằm trên một đường thẳng.}
\loigiai{
\begin{itemize}
    \item Đoạn thẳng đo được độ dài: Đúng.
    \item Đường thẳng kéo dài vô tận 2 phía: Đúng.
    \item 3 điểm bất kì có thể không thẳng hàng. Ý 3 Sai, Ý 4 Đúng.
\end{itemize}
\begin{itemize}
    \item a) \textbf{Đúng}.
    \item b) \textbf{Đúng}.
    \item c) \textbf{Sai}.
    \item d) \textbf{Đúng}.
\end{itemize}
}
\end{ex}

\begin{ex} %[2H5V25]
Cho ba điểm $A, B, C$ thẳng hàng theo thứ tự đó, $AB = 6\text{ cm}, BC = 8\text{ cm}$. Xét tính đúng sai:
\choiceTF
{\True Điểm $B$ nằm giữa hai điểm $A$ và $C$.}
{\True Đoạn thẳng $AC$ dài $14\text{ cm}$ ($6 + 8 = 14$).}
{Đoạn thẳng $AB$ dài hơn đoạn thẳng $BC$.}
{\True Đoạn thẳng $BC$ dài hơn đoạn thẳng $AB$ là $2\text{ cm}$ ($8 - 6 = 2$).}
\loigiai{
\begin{itemize}
    \item $B$ nằm giữa $A$ và $C$: Đúng.
    \item $AC = 6 + 8 = 14\text{ cm}$: Đúng.
    \item $AB (6\text{ cm}) < BC (8\text{ cm})$. Ý 3 Sai, Ý 4 Đúng.
\end{itemize}
\begin{itemize}
    \item a) \textbf{Đúng}.
    \item b) \textbf{Đúng}.
    \item c) \textbf{Sai}.
    \item d) \textbf{Đúng}.
\end{itemize}
}
\end{ex}

\textbf{Phần 3. Trả lời ngắn}
\setcounterpageref{ex}{0}

\begin{ex} %[2H5H25]
Cho đoạn thẳng $AB$ dài $12\text{ cm}$ và đoạn thẳng $CD$ dài $18\text{ cm}$. Tính tổng độ dài hai đoạn thẳng theo xăng-ti-mét.
\shortans{30}
\loigiai{
\textbf{Đáp số:} $30$\\
$12 + 18 = 30\text{ cm}$.
}
\end{ex}

\begin{ex} %[2H5H25]
Cho đoạn thẳng $MN$ dài $25\text{ cm}$, đoạn thẳng $PQ$ dài $14\text{ cm}$. Tính hiệu độ dài hai đoạn thẳng theo xăng-ti-mét.
\shortans{11}
\loigiai{
\textbf{Đáp số:} $11$\\
$25 - 14 = 11\text{ cm}$.
}
\end{ex}

\begin{ex} %[2H5V25]
Cho 3 điểm thẳng hàng $A, B, C$ với $B$ nằm giữa $A$ và $C$. Biết $AB = 8\text{ cm}, AC = 15\text{ cm}$. Tính độ dài $BC$ theo xăng-ti-mét.
\shortans{7}
\loigiai{
\textbf{Đáp số:} $7$\\
$BC = 15 - 8 = 7\text{ cm}$.
}
\end{ex}

\begin{ex} %[2H5V25]
Có 3 điểm không thẳng hàng, hỏi nối 3 điểm đó lại tạo thành bao nhiêu đoạn thẳng?
\shortans{3}
\loigiai{
\textbf{Đáp số:} $3$\\
Có 3 đoạn thẳng nối từng đôi một.
}
\end{ex}

\begin{ex} %[2H5V25]
Đoạn thẳng thứ nhất dài $9\text{ cm}$, đoạn thẳng thứ hai dài gấp đôi đoạn thứ nhất. Tính độ dài đoạn thứ hai theo xăng-ti-mét.
\shortans{18}
\loigiai{
\textbf{Đáp số:} $18$\\
$9 + 9 = 18\text{ cm}$.
}
\end{ex}

\begin{ex} %[2H5C25]
Trên một đường thẳng có 5 điểm phân biệt. Hỏi có tất cả bao nhiêu đoạn thẳng được tạo thành từ 5 điểm đó?
\shortans{10}
\loigiai{
\textbf{Đáp số:} $10$\\
Số đoạn thẳng: $4 + 3 + 2 + 1 = 10$ đoạn thẳng.
}
\end{ex}
